% id: 113-02
% book: 베이직쎈 확률과 통계 · 06 이항분포와 정규분포
% section: 유형 066 정규분포곡선의 성질
% figure: none
% answer: 
% answer_source: 
% variant_level: 0 (원본 전사)
% 이 파일은 build.mjs 가 items.json 에서 만든다. 고칠 때는 items.json 을 고친다.
\smash{\raisebox{1.5mm}{\dmkichul{베이직쎈 확률과 통계 113쪽 02번}}}
정규분포 $\mathrm{N}(m{,}\mkern3mu\mathopen{}\sigma^2)$을 따르는 확률변수 $X$의 확률밀도함수의 그래프의 성질로 옳은 것만을 \textbf{보기}에서 있는 대로 고른 것은? \bogi{ㄱ. $x=m$일 때 최댓값을 갖는다.\\ ㄴ. $\sigma$의 값이 일정할 때, $m$의 값이 작을수록 그래프의 대칭축은 오른쪽에 있다.\\ ㄷ. $\sigma$의 값이 같으면 그래프의 모양은 같다.\\ ㄹ. $\sigma$의 값이 클수록 그래프의 폭은 좁아진다.}
\dmchoicesauto{}{ㄱ, ㄴ}{ㄱ, ㄷ}{ㄴ, ㄷ}{ㄴ, ㄹ}{ㄷ, ㄹ}
